% STATUS: *draft* | reviewed | final
% PROMOTE: none | exemplars | quant-prep
% TAGS: arrays, hashing, dp, graphs, etc.
% DIFFICULTY: medium
% SOURCE: LeetCode 445

\ProblemTitle{Add Two Numbers II}
\ProblemMeta{Medium}{Linked Lists, Math, Stack}{LC-445}
\label{prob:lc-445-add-two-numbers-2}

\ProblemSummary
Given two non-empty linked lists representing two non-negative integers, where each node
stores a single digit and the first digits represent the \textbf{most significant} digit.
Add the two numbers and return the result as a \textbf{Linked List}.

 
\KeyIdea
TODO

% -----------------------------------------------------------------------------
% Optional: Author Thinking (excluded from exemplar builds)
% -----------------------------------------------------------------------------
\begin{Thinking}
My first thought upon reading this challenge is that, was to think back to LC-002 Add Two Numbers. 
Addition when the linked list was ordered by \textbf{least significant} digit, made addition simple. 
I could therefore reverse both input linked lists and apply the same algorithm for addition as per the previous challenge.
\end{Thinking}

% -----------------------------------------------------------------------------
% Algorithm
% -----------------------------------------------------------------------------
\Algorithm
\begin{CodeBlock}[style=pseudo,caption={Pseudocode}]{16}
    let l1, l2 be linked lists
    reverse l1, l2,

    iterate while l1, l2 or carry:
      maintaining pointers
      performing addition
      appending new node
      traversing

\end{CodeBlock}


% -----------------------------------------------------------------------------
% Correctness
% -----------------------------------------------------------------------------
\Correctness


% -----------------------------------------------------------------------------
% Complexity
% -----------------------------------------------------------------------------
\Complexity
Let $n$ and $m$ be the lengths of the two input linked lists.
Each iteration processes at most one node from each list, and the loop runs for
at most $\max(n, m) + 1$ iterations, yielding a time complexity of
$O(\max(n, m))$.

The algorithm uses $O(1)$ auxiliary space, as it maintains only a constant
number of pointers and a carry variable.
The output linked list requires $O(\max(n, m))$ space to store the result.


% -----------------------------------------------------------------------------
% Implementation
% -----------------------------------------------------------------------------
\bigskip
\noindent\textbf{C++ Implementation}

\lstinputlisting[
  style=cpp,
  caption={C++ Solution: LC-445 Add Two Numbers II}
]{../problems/06-linked-structures/medium/lc-445-add-two-numbers-2/solution.cpp}

\bigskip
\noindent\textbf{Python Implementation}

\lstinputlisting[
  language=Python,
  backgroundcolor=\color{codebg},
  basicstyle=\ttfamily\small,
  frame=single,
  breaklines=true,
  caption={Python Solution: LC-445 Add Two Numbers II}
]{../problems/06-linked-structures/medium/lc-445-add-two-numbers-2/solution.py}

\ProblemDivider
